% =====================================================================
% narrative-template.tex  --  the design contract for ONE paper
%
% WHAT THIS IS: a one-page, evidence-tracked story. It is NOT a draft of
% the paper. It mirrors the paper's REAL sections (Intro -> Methods ->
% Results -> Discussion), and every beat carries a readiness tag wired to
% the probe pipeline plus a small-font interrogation comment.
%
% HOW TO USE: copy this file to <paper>/0-lifecycle/3-narrative/3-narrative.tex,
% then replace every <...> placeholder. Delete beats you do not need; add
% beats the paper needs. Keep the readiness-tag macros and the legend.
% A filled, real-world exemplar lives at:
%   examples/ProjB-PhyTrait-OpioidRx/paper/Paper-Personality-Opioid-MedJournal/
%       0-lifecycle/3-narrative/3-narrative.tex
% =====================================================================
\documentclass[11pt]{article}
\usepackage[margin=1in]{geometry}
\usepackage{xcolor}
\usepackage{hyperref}

% --- readiness tags: a colored judgment on each narrative point ---
% Each beat gets exactly one tag. The tag couples the beat to its evidence
% status and routes [PENDING]/[GAP] beats to /haipipe-probe.
\definecolor{cReady}{rgb}{0.0,0.55,0.0}
\definecolor{cPend}{rgb}{0.85,0.50,0.0}
\definecolor{cLit}{rgb}{0.0,0.0,0.75}
\definecolor{cGap}{rgb}{0.80,0.0,0.0}
\definecolor{cInfer}{rgb}{0.45,0.10,0.55}
\newcommand{\rdy}{\textcolor{cReady}{\textbf{[READY]}}}   % evidence in hand: a confirmed probe or a run we trust
\newcommand{\pend}{\textcolor{cPend}{\textbf{[PENDING]}}} % data exists but a render/check/probe is still open
\newcommand{\lit}{\textcolor{cLit}{\textbf{[LIT]}}}       % rests on external literature; citation-audit pending
\newcommand{\gp}{\textcolor{cGap}{\textbf{[GAP]}}}        % no evidence yet; needs a probe
\newcommand{\infr}{\textcolor{cInfer}{\textbf{[INFER]}}}  % an inference: grounded in evidence, one reasoned step beyond, never measured
% \rev = small-font INTERNAL interrogation comment attached BELOW a beat.
% #1 = verb $\cdot$ role  (verb in {keep, add, demoted, cut, added by author};
%      role in {stakes, validity, contribution, guardrail, safety, defense,
%      mechanism, grounded opener, no-blame anchor, so-what, ...})
% #2 = one sharp sentence: why this beat is here / what breaks without it.
\newcommand{\rev}[2]{\\{\footnotesize\color{gray}$\hookrightarrow$~\textbf{#1}~#2}}

% \fb = an EXTERNAL reviewer's comment, threaded onto the beat it concerns
% (post + comments model). Use when a co-author / advisor / referee reviews the
% paper and you want each comment tracked at the point where the change happens.
% Distinct maroon, so it reads apart from the gray internal \rev.
% #1 = reviewer name        (e.g. Ritu)
% #2 = status               (done | part | open)
% #3 = their feedback, VERBATIM (their words; do not paraphrase or compress;
%      only escape LaTeX specials % & _ $ # { } so it compiles)
% #4 = how we addressed it  (OUR words; SHORT PLAIN sentences, one idea each)
% Order on a beat that has both: beat text -> \rev (internal) -> \fb (external).
% Multiple reviewers: each gets its own footer "External review (<name>, <date>)"
% line; their \fb lines coexist on a beat (each carries its own name).
% A comment with no single home beat (e.g. "see comments in Overleaf") stays in
% the footer ledger, not on a beat.
\definecolor{cFb}{rgb}{0.55,0.0,0.15}
\newcommand{\fb}[4]{\\{\footnotesize\color{cFb}$\hookrightarrow$~\textbf{#1}~[\textbf{#2}]~\emph{``#3''}~$\Rightarrow$~#4}}

\title{Narrative: The Paper as One Story}
\date{<YYYY-MM-DD>}
\begin{document}
\maketitle

% --- The legend is printed so a reader knows what each tag means. Keep verbatim. ---
\noindent\footnotesize\textbf{Readiness legend (per point):}
\rdy~evidence in hand (a confirmed probe or a run we trust) $\cdot$
\pend~data exists but a render/check/probe is still open $\cdot$
\infr~an inference: grounded in the evidence, reaching one reasoned step beyond, not measured (no probe will confirm it) $\cdot$
\lit~rests on outside literature, citation-audit pending $\cdot$
\gp~no evidence yet, needs a probe.\normalsize

% --- The SPINE is the whole paper in one breath: an arrow chain from the ---
% --- problem to the so-what. Everything below must serve this line. ---
\medskip
\noindent\textbf{Spine (throughline):} <the problem / stakes> $\rightarrow$ <the move or signal this paper introduces> $\rightarrow$ <the core finding> $\rightarrow$ <the so-what / what it is and is not>.

% Format reminder: each section = a short story PARAGRAPH (the narrative flow,
% draft-quality, plain language, Pn.Sm tagged) + "Key points to cover" BULLETS
% behind it. Each bullet = readiness tag + bold label + 1-3 sentences + a \rev
% interrogation comment. The whole walks Intro -> Methods -> Results ->
% Discussion. Adapt section subtitles, Flow lines, and beats to the venue.
%
% COMMENT STYLE (both \rev and \fb): write SHORT PLAIN sentences, one idea each.
% No run-on lines chained by semicolons, no stacked parentheticals. In \footnotesize
% a long compound line is unreadable; 2-4 short sentences scan in one pass. Lead
% with the action, then the open item. (Same readability rule as the pitch.)
%
% EXTERNAL REVIEW: when a named reviewer comments on the paper, thread each comment
% onto the beat it concerns with \fb (see the example under the first Intro beat).
% Keep their feedback verbatim; keep your resolution short. \fb is for external
% named reviewers; \rev is the internal interrogation pass.

% =====================================================================
\section*{Introduction: <what is known, the gap, and the bet>}

\noindent\textbf{Flow:} <known fact> $\rightarrow$ <what exists at scale> $\rightarrow$ <the precise unmet gap> $\rightarrow$ <so we test X>.

%% ---- P1.S1 ----
<One grounded sentence: the broad fact that gives the paper stakes.>
%
%% ---- P1.S2 ----
<One sentence: the data or signal that makes the question answerable.>
%
%% ---- P1.S3 ----
<One sentence: the precise gap, then the question we ask.>

Key points to cover:
\begin{enumerate}
\item \lit~\textbf{<Stakes>:} <the external fact that makes this matter, with a citation>. \rev{keep $\cdot$ stakes}{<why this beat earns its place; the exact link to cite so the stakes do not read as borrowed>.}
% EXAMPLE of an external reviewer comment threaded onto a beat (delete if unused):
\fb{<Reviewer>}{part}{<their comment, verbatim>}{<how we addressed it, in short plain sentences. One idea each.>}
\item \rdy~\textbf{<The gap>:} <draw the white space precisely: what prior work did and did not do>. \rev{keep $\cdot$ contribution setup}{<why this is the strongest Intro beat; what to scope so the contribution is not overclaimed>.}
\item \rdy~\textbf{<Our contribution>:} <the one-sentence reason this paper exists; name any enabling method as enabler, not contribution>. \rev{keep $\cdot$ contribution}{<keep the scope tight; keep enabler-vs-contribution clear here, not floated to Methods>.}
\end{enumerate}

% =====================================================================
\section*{Methods: <how we test it>}

\noindent\textbf{Flow:} <the cohort/setting> $\rightarrow$ <the build> $\rightarrow$ <the exposure/measure> $\rightarrow$ <the outcomes> $\rightarrow$ <the estimator and its honesty>.

%% ---- P2.S1 ----
<One sentence: setting + the measure, flagged as enabler if it is one.>
%
%% ---- P2.S2 ----
<One sentence: how the analytic sample is built, so it is reproducible.>
%
%% ---- P2.S3 ----
<One sentence: what we relate to what, and the honesty guardrail (e.g. associational).>

Key points to cover:
\begin{enumerate}
\item \rdy~\textbf{<Setting>:} <population and N; define the unit of analysis in one clause>. \rev{keep $\cdot$ validity}{<frames the denominator before any estimate; define the unit or the population is ambiguous>.}
\item \pend~\textbf{<Cohort construction>:} <the build rules, pinned to code; note any number still server-only or unbuilt>. \rev{keep $\cdot$ validity}{<the reproducibility spine; until the pending number lands it is a liability, not a strength>.}
\item \rdy~\textbf{<Exposure / measure>:} <the key variable; inoculate against the obvious construct-validity attack>. \rev{keep $\cdot$ validity}{<where the construct-validity attack is met; carry the enabler flag explicitly if it applies>.}
\item \rdy~\textbf{<Outcomes>:} <the dependent measures; define thresholds here>. \rev{keep $\cdot$ validity}{<ties the test to what makes the result matter; define thresholds where they are first used>.}
\item \rdy~\textbf{<Design>:} <the estimator + key controls + the honesty guardrail>. \rev{keep $\cdot$ guardrail}{<carries the honesty (e.g. associational); pre-specify primary contrasts or they read as cherry-picked>.}
\end{enumerate}

% =====================================================================
\section*{Results: <what the data reveal>}

\noindent\textbf{Flow:} <who is in the sample> $\rightarrow$ <the main result> $\rightarrow$ <where it concentrates> $\rightarrow$ <the high-risk tail> $\rightarrow$ <it survives the obvious confound> $\rightarrow$ <it holds under the robustness check>.

%% ---- P3.S1 ----
<One sentence: describe the cohort / Table 1.>
%
%% ---- P3.S2 ----
<One sentence: the headline finding, surprise kept sharp.>
%
%% ---- P3.S3 ----
<One sentence: where it amplifies and what shape it takes.>

Key points to cover:
\begin{enumerate}
\item \pend~\textbf{<Cohort (Table 1)>:} <who is in the sample and how key variables distribute>. \rev{keep $\cdot$ stakes}{<the section cannot open without it; mark PENDING until the table is built>.}
\item \rdy~\textbf{<C1, main>:} <the headline result in one line>. \rev{keep $\cdot$ contribution}{<the headline; keep any surprise sharp>.}
\item \rdy~\textbf{<C2, where it concentrates (safety/impact)>:} <the result that carries the venue fit>. \rev{keep $\cdot$ contribution}{<present with the pre-specified anchor and one reconciled number set, or it dies on multiplicity>.}
\item \rdy~\textbf{<C3, the tail>:} <the result that converts an average into a consequential claim>. \rev{keep $\cdot$ safety}{<this is the beat the venue rewards; lead it over shape detail>.}
\item \pend~\textbf{<Robustness>:} <the check that defends the main result against the obvious "is it just outliers" reflex>. \rev{add $\cdot$ defense}{<describe the check, not its outcome; once run it flips to READY>.}
\item \gp~\textbf{<Mechanism split>} (parked, not in Results): <the decomposition that needs a probe before it can be claimed>. \rev{cut from Results $\cdot$ mechanism}{<why it is invalid to assert now; it returns as a Discussion mechanism beat after the probe>.}
\end{enumerate}

% =====================================================================
\section*{Discussion: <what it means>}

\noindent\textbf{Flow:} <finding restated> $\rightarrow$ <most plausible explanation, no blame> $\rightarrow$ <what it means in practice> $\rightarrow$ <what the result is good for> $\rightarrow$ <what it cannot claim>.

%% ---- P4.S1 ----
<One sentence: restate the principal finding, no new numbers.>
%
%% ---- P4.S2 ----
<One sentence: the most plausible reading, framed as plausible not proven.>
%
%% ---- P4.S3 ----
<One sentence: the practical / bedside / field implication.>
%
%% ---- P4.S4 ----
<One sentence: the so-what plus the honest scope of the evidence.>

Key points to cover:
\begin{enumerate}
\item \rdy~\textbf{<Principal finding>:} <the restated headline a Discussion opens on; pure summary, no new numbers>. \rev{add $\cdot$ grounded opener}{<the grounded anchor before any interpretation>.}
\item \infr~\textbf{<Mechanism>:} <the most plausible explanation; cite the supporting literature here only>. \rev{keep $\cdot$ guardrail}{<phrase as "most plausible reading," never as established; this beat owns the relocated mechanism literature>.}
\item \infr~\textbf{<Implication / so-what>:} <what a reader should take away>. \rev{keep $\cdot$ so-what}{<the so-what must be in the sentence itself, not implied>.}
\item \rdy~\textbf{<Limitations>:} <the honest list: design limits, construct limits, confounds the main claim must pre-empt>. \rev{keep $\cdot$ guardrail}{<mandatory; the limitation that pre-empts the strongest confounding attack goes here>.}
\item \pend~\textbf{<Multiple comparisons / other reviewer rejection vector>:} <the honest accounting of the most predictable rejection vector at this venue>. \rev{add $\cdot$ guardrail}{<pairing the honest burden with the pre-specified anchor is what lets the peak claim survive it>.}
\end{enumerate}

% =====================================================================
% FOOTER: the cross-cutting ledger. Reviewer-flagged gaps thread back into
% the beats above; the Arc states how the spine's peak claim is defended;
% Awaiting review names beats that still need an independent interrogation pass.
% External reviewer comments (\fb) live ON their beats above; the footer only
% carries comments with no single home beat (e.g. "see comments in Overleaf").
% =====================================================================
\bigskip
\noindent{\footnotesize\color{gray}\textbf{Reviewer-flagged gaps:} <list each known reviewer concern and where it is now threaded (which section beat), or mark it Remaining and route it to a probe/computation>.
\textbf{Arc:} <one line: after the demotions/parks/folds above, what does each section land on, and how is the spine's peak claim defended>.
\textbf{Awaiting review:} <any beat authored since the last interrogation pass that still needs an independent keep/move/demote/cut verdict>.
\textbf{External review (<name>, <date>):} <one line pointing to the threaded \fb comments above + any comment that has no home beat (e.g. "see Overleaf"); source file path>.}

\end{document}
